As operações orbitais de Rendezvous and Docking or Berthing (RVD/B) é um procesos através do qual duas espaçonaves se aproximam uma da outra, envolvendo procedimentos de alinhamento e sincronização de movimentos para realizar com segurança operações de acoplamento em portas de docking ou realizar operações de ancoramento na nave alvo via manipulador robótico que pode estar no alvo na nave visitante. O RVD/B é essencial para tarefas como montar grandes naves espaciais, consertar plataformas e estações orbitais e trocar pessoal entre estações, consertar naves espaciais no espaço e recuperar naves espaciais, incluindo sua captura e retorno ao mundo, ou seja, é essencial para prestação de serviços em órbita (também conhecido com OOS, on orbit services), essas ideias apresentadas também permitem o uso de veículos em órbita com espaçonaves retornando à Lua e à Terra \ \citep{fehse2003}.

O primeiro rendezvous and docking (RVD) entre duas espaçonaves ocorreu em 1966 \citep{nasa1stDocking}, quando Neil Armstrong e Dave Scott realizaram manualmente um \textit{rendezvous} (aproximação) em uma nave Gemini e, em seguida, acoplaram com um veículo-alvo não tripulado, denominado Agena. Já o primeiro RVD automático aconteceu em 1967, com o acoplamento das espaçonaves soviéticas Cosmos 186 e 188 \citep{fehse2003}. Desde então, essas operações têm sido amplamente aplicadas em programas espaciais da Rússia (antiga União Soviética) e dos Estados Unidos, como os programas Apollo (1968–1972) e Skylab (1973–1974), além dos projetos soviéticos Salyut e Mir (1971–1999), que envolveram o uso de naves tripuladas Soyuz e não tripuladas Progress \citep{fehse2003}.

Na Europa Ocidental, a tecnologia de RVD/B foi desenvolvida pela Agência Espacial Europeia (ESA) desde a década de 1980. Inicialmente classificada como "tecnologia habilitadora", essas técnicas foram implementadas no projeto Columbus Man-Tended Free-Flyer (MTFF), planejado para acoplar-se à estação espacial americana Freedom \citep{goelz1992MtffFreedom} e à espaçonave europeia Hermes. Após o cancelamento dos projetos MTFF e Hermes, o Veículo de Transferência Automatizado (ATV) tornou-se a principal contribuição europeia para o Programa da Estação Espacial Internacional (ISS) \citep{ellwood2007ATV2ISS}. O ATV foi concebido para realizar tarefas essenciais de abastecimento e manutenção na ISS, um passo importante na integração espacial europeia \citep{fehse2003}.

Os veículos que operaram RVD/B na ISS incluem o Ônibus Espacial dos EUA, a Soyuz russa (tripulada), a espaçonave Progress (não tripulada), o ATV europeu e a aeronave de transporte japonesa HTV. Outras aplicações do RVD/B, como missões de ônibus espaciais, recuperação de telescópios espaciais e missões de retorno lunar e planetário, ainda continuam a ser exploradas.

\subsection{Processo de RVD/B}

O processo de rendezvous and docking or berthing envolve uma série de manobras e trajetórias controladas, que trazem gradualmente o veículo \textit{chaser} (perseguidor) para a proximidade do veículo \textit{target} (alvo). Esse processo exige um posicionamento preciso em termos de posição, velocidade e atitude, atendendo aos critérios técnicos para realização da manobra.
No acoplamento (docking), o sistema de controle do \textit{chaser} ajusta os parâmetros necessários para conectar-se diretamente às interfaces do \textit{chaser}.
Já no berthing (atracação), o \textit{chaser} é levado a um ponto de encontro, onde um manipulador finaliza o acoplamento \citep{fehse2003}.

A complexidade dessas operações resulta das diversas condições e restrições envolvidas, como a necessidade de alinhar o perseguidor e o alvo no mesmo plano orbital e ajustar as trajetórias em resposta a perturbações orbitais.
As estratégias de lançamento e faseamento são essenciais para garantir a proximidade dos veículos no momento adequado. Além disso, as trajetórias devem ser continuamente atualizadas para compensar alterações nos parâmetros orbitais causadas por fatores externos \citep{fehse2003}.

\subsection{Programa espacial Apollo}

O programa Apollo marcou uma transição significativa na arte de pilotar veículos espaciais, introduzindo novas dimensões operacionais que ampliaram as capacidades tanto da tripulação quanto dos sistemas embarcados. Três áreas principais foram destacadas: as operações de voo (Flight Operations), o papel da tripulação (Flight Crew's Role) e as comunicações homem-máquina (Man-Machine Communications). Essas dimensões refletiam o alto nível de integração necessário entre os controladores em terra e os sistemas a bordo, bem como a complexidade das funções desempenhadas pela tripulação \citep{nevins1971}.

Para a tripulação, os aspectos principais incluíam a amplitude das tarefas gerais que precisavam ser realizadas e a necessidade de supervisionar sistemas críticos de voo. Isso envolvia a execução de tarefas em diferentes níveis de complexidade, como o gerenciamento de sistemas e a navegação precisa da espaçonave. Adicionalmente, a interação com os equipamentos da nave evoluiu de uma atuação direta nos controles para um nível mais avançado de comunicação funcional com os sistemas \citep{nevins1971}.

Um exemplo dessa transição foi o sistema de controle, navegação e orientação desenvolvido especificamente para as naves Apollo. Esse sistema foi projetado para fornecer à tripulação um ambiente operacional autônomo, permitindo que os astronautas navegassem e guiassem o veículo sem dependência contínua do suporte em terra. Essa capacidade foi decisiva para o sucesso das missões, garantindo que as operações pudessem ser conduzidas mesmo em situações imprevistas ou fora do alcance da comunicação direta com a Terra \citep{nevins1971}.

Para que a operação autônoma pudesse acontecer, \citep{nevins1971} explica quais requisitos foram estabelecidos e que deveriam ocorrer para a missão Apollo: (i) o sistema deveria ser capaz de completar a missão sem auxílio ou ajuda da estação em solo; (ii) o sistema empregaria a participação humana sempre que isso pudesse simplificar ou melhorar a operação em comparação às sequências automáticas das funções exigidas; (iii) o sistema deveria fornecer interfaces adequadas ao piloto, incluindo displays (ou telas e informações adequadas), e métodos de controle do sistema de orientação; (iv) o sistema deveria ser projetado de modo que um único membro da tripulação fosse capaz de executar todas as funções necessárias para garantir retorno seguro à Terra a partir de qualquer ponto da missão.

Como resultado, o Apollo tornou-se o primeiro programa tripulado dos Estados Unidos a integrar sensores e sistemas avançados de processamento de dados a bordo, capacitando a tripulação a desempenhar um papel ativo e decisivo nas operações de voo. Essa abordagem não apenas ampliou a eficiência das missões, mas também estabeleceu novos padrões para a interação entre o homem e a máquina em voos espaciais \citep{nevins1971}.

Para estabelecer um melhor entendimento sobre a interação entre o homem e a máquina, é necessário entender que a comunicação entre os tripulantes e os sistemas de bordo na missão Apollo era realizado desde a execução direta de tarefas, por meio de atuadores manuais, até um nível funcional de comunicação via uma linguagem computacional de nível superior \citep{nevins1971}, tal linguagem permitia que os astronautas controlassem grupos de tarefas em vez de tarefas individuais, isso facilitava o gerenciamento de operações complexas como, por exemplo, a rotação automática da atitude da espaçonave ou até a integração de diversas ações necessárias para o pouso na Lua ou também, como exemplo, a reentrada na atmosfera terrestre.

Para se ter uma ideia, durante a fase de descida motorizada no pouso lunar, o controle da atitude da espaçonave envolvia quatro modos no sistema primário: (i) manobra automática para um posicionamento predefinido; (ii) manutenção da atitude com parâmetros ajustáveis; (iii) manobra manual com controlador dos três eixos; e (iv) manobras de pequena amplitude com impulsos temporizados. O sistema de backup, que possuía caminhos paralelos, também operava com quatro modos, ativados manualmente e de forma similar.

O papel da tripulação nas missões Apollo foi fundamental para a navegação e o controle da espaçonave, abrangendo desde a supervisão de sistemas automáticos até a realização de tarefas específicas de monitoramento, decisão e operação. Esse papel pode ser categorizado em diferentes funções principais, cada uma delas exigindo um alto nível de integração entre o ser humano e a máquina \citep{nevins1971}.

Entre as tarefas gerais realizadas pela tripulação, destacam-se:

a) Monitoramento e tomada de decisões associadas aos processos de navegação e orientação, incluindo a análise de dados dos sensores de navegação (como rastreamento de alvos, visual e por radar) e a atualização de vetores de estado. A tripulação comparava os dados a bordo com os dados de rastreamento em terra e gráficos de backup para garantir a precisão das operações.

b) Sequenciamento e inicialização dos sistemas primários de orientação, navegação e sensores, bem como dos sistemas de propulsão e temporização.

c) Inicialização e sequenciamento de sistemas de backup, garantindo redundância e confiabilidade em caso de falhas dos sistemas principais.

d) Realização de tarefas de reconhecimento de padrões, como o rastreamento óptico do módulo lunar durante operações de rendezvous, aquisição e identificação de estrelas, além do alinhamento geométrico com o horizonte visual (Terra ou Lua) para navegação cislunar.

Essas tarefas eram essenciais para o sucesso das missões, pois a tecnologia disponível na época ainda não permitia a total automação dessas operações. Embora os avanços tecnológicos gradualmente tenham melhorado a capacidade e a confiabilidade dos sistemas, muitos desses processos continuavam exigindo a intervenção humana direta durante o programa Apollo. 

A complexidade das missões Apollo, dependia de um design que integrava profundamente homem e máquina. Essa integração não apenas proporcionou flexibilidade nas operações, mas também garantiu a confiabilidade necessária para enfrentar os desafios únicos das missões espaciais. Assim, o papel da tripulação foi estrategicamente projetado para maximizar o desempenho, utilizando a combinação de habilidades humanas e automação para alcançar resultados que não seriam possíveis apenas com a tecnologia disponível \citep{nevins1971}.

À medida que a tecnologia avançava em capacidade e confiabilidade, muitos dos sistemas utilizados no programa Apollo começaram a permitir maior automação, reduzindo a necessidade de intervenção constante da tripulação, no entanto, o papel humano ainda continuou sendo essencial devido às limitações tecnológicas da época \citep{nevins1971}.

As tarefas de supervisão desempenhadas pela tripulação eram especialmente críticas, considerando a confiabilidade limitada dos componentes básicos, aliada às restrições de peso, volume e potência. Essas condições resultavam em projetos de sistemas onde uma única falha poderia incapacitar funcionalidades importantes. Para mitigar esses riscos, era necessário incorporar sistemas redundantes. A função mais importante da tripulação, especialmente em condições dinâmicas, era monitorar tanto os sistemas primários quanto seus sistemas de backup associados, garantindo transições suaves e rápidas entre eles, quando necessário. Essa função exigia atenção contínua e compartilhada entre vários níveis de sistemas redundantes \citep{nevins1971}.

Além disso, era essencial que a transição para modos de backup ocorresse de maneira gradual, ao invés de abrupta, portanto a degradação do sistema acontecia de forma lenta ao invés de instantânea, aumentando assim, a robustez operacional. No entanto, nos cinco a dez anos seguintes ao programa Apollo, ainda era improvável que a flexibilidade e a capacidade de decisão humanas fossem totalmente substituídas por sistemas automatizados. Durante esse período, os sistemas precisariam ser configurados para manter o envolvimento humano, embora de forma reduzida e voltada principalmente para funções de supervisão ou monitoramento \citep{nevins1971}.

Com o passar do tempo e o avanço tecnológico, esperava-se que o papel da tripulação se tornasse mais administrativo e supervisionado, aliviando significativamente a carga operacional. Essa evolução permitiria que os sistemas automáticos assumissem a maior parte das funções críticas, enquanto o papel humano seria focado na supervisão geral e na intervenção em situações excepcionais \citep{nevins1971}.

A comunicação entre a tripulação e os equipamentos a bordo das espaçonaves Apollo desempenhou um papel crucial no sucesso das missões, abrangendo desde a execução direta de tarefas específicas até níveis mais elevados de interação mediados por sistemas computacionais avançados. Essa interação variava entre a atuação direta nos controladores manuais para a execução de tarefas específicas e o uso de uma linguagem computacional de alto nível, que permitia à tripulação controlar grupos de tarefas de forma integrada \citep{nevins1971}.

A linguagem computacional implementada nos sistemas Apollo possibilitou a realização de tarefas que iam desde a rotação automática da atitude da espaçonave até a integração de operações complexas, como a descida controlada durante o pouso lunar ou a reentrada na atmosfera terrestre. Essa inovação foi uma contribuição significativa para a evolução na arte de pilotar veículos espaciais, permitindo maior eficiência e precisão nas operações \citep{nevins1971}.

Embora a linguagem computacional dos sistemas Apollo fosse avançada para a época, apresentava limitações em comparação às tecnologias atuais, como a ausência de interfaces gráficas e formatos de entrada e saída de dados mais flexíveis. No entanto, seu desenvolvimento foi essencial para facilitar a interação homem-máquina, proporcionando à tripulação ferramentas que melhoraram a gestão das operações e aumentaram a confiabilidade dos sistemas \citep{nevins1971}.

Além disso, o sistema de gerenciamento de voo da Apollo integrava diversos aspectos de navegação, orientação e controle, permitindo que a tripulação supervisionasse e controlasse operações críticas com maior eficiência. Essa abordagem foi aplicada em várias fases das missões, incluindo a descida controlada para o pouso lunar e os modos de controle de atitude rotacional do módulo de comando. Essas funcionalidades não apenas aumentaram a segurança das missões, mas também representaram um marco no desenvolvimento de sistemas avançados de comunicação homem-máquina \citep{nevins1971}.

